\chapter{결론}\label{ch:conclusion}

\section{지금까지의 결과}

이 프로젝트의 최종 목표는 관람객이 직접 조작할 수 있는 기계식 봄베이고, 전시의 구성은 기계식, 릴레이식,
전자식(FPGA)의 세 대를 나란히 두어 계산 기계의 발전을 보이는 것이다. 지금까지 끝난 것은 세 대 모두가
내야 할 답을 정의하는 기준 모델과, 전자식 구현의 기준이 되는 회로다. 각 기계의 스크램블러가 올바른지,
올바른 위치에서 정지하는지는 이 기준과 비교해 판정한다.

세 기계의 현재 상태는 서로 다르다(표~\ref{tab:status}). 아직 어느 기계도 실물로 동작한 적이 없다.

\begin{table}[ht]
\centering
\caption{세 기계의 현재 상태(2026년 11월).}
\label{tab:status}
\small
\begin{tabularx}{\linewidth}{>{\raggedright\arraybackslash}p{22mm}>{\raggedright\arraybackslash}X>{\raggedright\arraybackslash}X}
\toprule
기계 & 끝난 것 & 남은 것 \\
\midrule
전자식 (FPGA) & 엔진의 설계와 시뮬레이션 검증. Nexys Video용 배치와 배선, 엔진 두 대가 50~MHz에서 타이밍을
  만족(\ref{sec:fpga-results}절). 호스트 소프트웨어 & 보드의 확보, 보드에서의 동작과 속도의 측정, 표시 장치 \\
기어식 & 개념 설계와 상세 입체 모델(\ref{sec:hw-gear-cad}절). 로터 I에서 V까지 드럼 기판의 설계와
  검증(\ref{sec:hw-drum-pcb}절). 접점 보드, 열 단자, 반사판 카드, 제어 보드의 회로도와 배치(\ref{sec:hw-boards}절)
  & 기판의 배선, 기판과 접점의 시제품, 자리올림 기구, 스크램블러 한 유닛 \\
릴레이식 & 개념 설계. 시프터와 계수기의 기준 모델을 코어에 넣어 모든 위치에서 검증(\ref{sec:hw-relay-model}절).
  시프터, 배선 카드, 계수기, 순서기, 위치 표시, 오프셋 다이얼의 회로도와 배치(\ref{sec:hw-boards}절)
  & 기판의 배선, 시프터 보드의 시제품, 스크램블러 한 유닛 \\
공통 & 설계 요구사항과 인터페이스의 정의(결정 0003), 전체 구동용 메뉴 카드(\ref{sec:hw-cards}절). 테스트 레지스터,
  대각선 보드 다섯 장, 패치 패널, 단위 시험기의 회로도와 배치 & 구간용 카드, 기판의 배선과 제작, 패널의 구멍 배치,
  실측 \\
\bottomrule
\end{tabularx}
\end{table}

소프트웨어 단계에서 확인한 것은 다음과 같다.

\begin{itemize}
\item 실제 기계의 구성요소를 일대일로 옮긴 모델이 크립에서 출발해 로터 순서, 코어 위치, 스테커, 평문까지
복원한다. 3로터와 4로터 모두에서, 그리고 1945년의 실제 전문에서 확인했다.
\item 거짓 정지의 수가 튜링의 추정표와 같은 비율로 변한다. 루프 하나가 정지를 약 $1/26$로 줄이고,
대각선 보드가 없으면 루프가 적은 메뉴는 쓸 수 없다.
\item 전류의 확산은 선을 살피는 순서와 무관한 하나의 집합을 구하는 일이다. 이 성질이 한 가닥씩 살피는
소프트웨어, 클럭마다 한 걸음씩 퍼뜨리는 회로, 전류가 저절로 퍼지는 실물이 같은 답을 내는 근거다.
기준 RTL이 Rust 코어와 676비트 단위로 일치하는 것으로 앞의 둘을 확인했다.
\end{itemize}

구현 과정에서 되풀이해 나타난 생각은 ``바뀌는 것만 다시 계산한다''였다. 실제 봄베의 느린 드럼,
소프트웨어의 스크램블러 표, FPGA 엔진의 표 메모리가 모두 같은 구조다.

\begin{todobox}{최종 결론}
FPGA와 실물 단계가 끝난 뒤 다음을 쓴다. 같은 메뉴에서 세 기계(기계식, 릴레이식, FPGA)의 정지가 기준
모델과 일치한 결과, 위치당 시간으로 본 기계에서 전기기계, 전자로의 속도 변화와 규모의 비교, 전시·실습
요구사항의 충족 여부, 실제 봄베의 설계 결정(드럼의 속도 배치, 자리올림 방식, 정지 감지)을 직접 만들어 보며 이해하게 된
점, 남은 과제(지시 드럼의 역사적 표기, 부가 장치, 확인하지 못한 자료).
\end{todobox}
